\section{Síntesis y ampliación}

\subsection{Puntos clave de esta clase}

\begin{enumerate}
    \item \textbf{Las imágenes son números}: la computadora representa una imagen mediante una matriz de números de $H \times W \times C$, donde $C$ es el número de canales (imagen en escala de grises $C=1$, imagen a color $C=3$).

    \item \textbf{La extracción de características es el núcleo del reconocimiento visual}: para reconocer un objeto es necesario extraer características que distingan entre categorías distintas. El diseño manual de características enfrenta el dilema de la recursividad y la variabilidad; el aprendizaje profundo resuelve este problema aprendiendo automáticamente una jerarquía de características a partir de los datos.

    \item \textbf{Las redes totalmente conectadas no son adecuadas para imágenes}: la operación de aplanado hace perder la información espacial y el número de parámetros resulta excesivo.

    \item \textbf{Las tres ventajas principales de la operación de convolución}:
    \begin{itemize}
        \item Conectividad local: preserva la estructura espacial
        \item Compartición de parámetros: reduce drásticamente el número de parámetros
        \item Equivariancia a la traslación: el mismo patrón se detecta sin importar en qué posición de la imagen aparezca
    \end{itemize}

    \item \textbf{Arquitectura de la CNN}: el apilamiento de capas de convolución + ReLU + capas de agrupamiento constituye el módulo de aprendizaje de características, y las capas totalmente conectadas + Softmax constituyen el módulo de clasificación.

    \item \textbf{La CNN es una arquitectura de visión de propósito general}: el mismo marco de aprendizaje de características puede usarse en múltiples tareas, como clasificación, detección de objetos y segmentación semántica.
\end{enumerate}

\subsection{Relación con las dos clases anteriores}

\begin{table}[H]
\centering
\begin{tabular}{p{3cm}p{5cm}p{5cm}}
\toprule
\textbf{Aspecto} & \textbf{Lecture 1 (red totalmente conectada)} & \textbf{Lecture 3 (CNN)} \\
\midrule
Forma de la entrada & Vector de características unidimensional & Tensor de imagen bidimensional/tridimensional \\
Forma de conexión & Totalmente conectada (todas las entradas a todas las neuronas) & Conectividad local (cada neurona solo observa una región local) \\
Compartición de parámetros & Ninguna & El filtro se comparte espacialmente \\
Información espacial & No se preserva & Se preserva y se aprovecha \\
Operación central & Multiplicación de matrices & Operación de convolución \\
\bottomrule
\end{tabular}
\caption{Comparación entre la red totalmente conectada y la CNN}
\end{table}

\begin{importantbox}{Filosofía central: dejar que la estructura de los datos guíe el diseño de la arquitectura}
La filosofía de diseño de la CNN consiste en \textbf{aprovechar la estructura intrínseca de los datos para diseñar la arquitectura de la red}. Las imágenes tienen localidad espacial: la relación entre píxeles vecinos es más estrecha que entre píxeles distantes; las imágenes tienen invariancia a la traslación: el mismo patrón puede aparecer en cualquier posición. La operación de convolución aprovecha perfectamente estos dos conocimientos previos. Esta filosofía se repite una y otra vez en el aprendizaje profundo: la RNN aprovecha la estructura temporal de las secuencias, y el mecanismo de atención del Transformer aprovecha la estructura de relaciones entre elementos.
\end{importantbox}

\subsection{Lecturas adicionales}

\begin{itemize}
    \item Sitio oficial del curso MIT 6.S191: \url{http://introtodeeplearning.com}
    \item LeCun et al. (1998), "Gradient-Based Learning Applied to Document Recognition" — obra fundacional de la CNN
    \item Krizhevsky et al. (2012), "ImageNet Classification with Deep Convolutional Neural Networks" (AlexNet) — hito del aprendizaje profundo en el campo de la visión por computadora
    \item Lee et al. (2009), ICML — artículo original de la visualización de la jerarquía de características presentada en esta clase
    \item McKinney et al. (2020), Nature — estudio en el que la CNN supera a los médicos humanos en la detección de cáncer de mama
    \item Esteva et al. (2017), Nature — estudio sobre el uso de la CNN para el diagnóstico de cáncer de piel
    \item Rohrer, "How do CNNs work?" — referencia del caso de detección de la X presentado en esta clase
    \item Karn, "An Intuitive Explanation of CNNs" — referencia de la animación de la operación de convolución presentada en esta clase
\end{itemize}

%% --- Fin del contenido --- %%

\end{document}
